\documentclass[a4paper,11pt]{ltjsarticle}
\usepackage[haranoaji]{luatexja-preset}
\usepackage{amsmath,amssymb}
\usepackage[top=12mm,bottom=10mm,left=14mm,right=14mm]{geometry}
\usepackage{array}
\usepackage{tikz}
\usetikzlibrary{calc}
\pagestyle{empty}
\setlength{\parindent}{0pt}
\newlength{\qcol}\setlength{\qcol}{0.47\textwidth}
\newlength{\qgap}
\newlength{\anscolw}\setlength{\anscolw}{0.30\textwidth}
\newlength{\ansh}
\newcommand{\Q}[2]{\parbox[t]{\qcol}{\raggedright (#1)\hspace{0.6em}$\displaystyle #2$}}
\newcommand{\QR}[4]{\Q{#1}{#2}\hfill\Q{#3}{#4}\par\vspace{\qgap}}
\newcommand{\anscell}[1]{\rule[-\ansdrop]{0pt}{\ansh}(#1)}
\newlength{\ansdrop}

\begin{document}
{\large\bfseries 計算問題　20}\hfill 名前(\underline{\hspace{32mm}})\par\vspace{3mm}
■（1）〜（16）の計算をしなさい。（17）〜（21）は方程式を解きなさい。\par\vspace{3mm}
\setlength{\qgap}{5.4mm}
\QR{1}{98x\div(-7)}{2}{(-8)\times9a}
\QR{3}{(-13)-(-8)}{4}{\dfrac{1}{3}(9x-6)+\dfrac{1}{2}(8x+10)}
\QR{5}{(-51)\div(-3)}{6}{-16\times(-13)}
\QR{7}{(-13)+(-18)}{8}{\dfrac{3}{4}(4x-8)-\dfrac{4}{5}(5x+10)}
\QR{9}{-3+2^4+(-3)^2}{10}{-0.7-(-2.4)-3.3}
\QR{11}{\dfrac{2}{3}\div2+2\times\left(-\dfrac{5}{6}\right)}{12}{7\times\{3-(-9+1)\times5\}}
\QR{13}{(2x+7)+(3x-8)}{14}{(-8)+(+12)-(-10)}
\QR{15}{2-(-5)\div\left(-\dfrac{5}{4}\right)}{16}{\dfrac{x-9}{2}\times8}
\QR{17}{3x-5=5x-7}{18}{-3x+2=11}
\QR{19}{\dfrac{2}{3}x+1=\dfrac{5}{6}x-1}{20}{0.9x-0.4=5}
\vspace{1mm}
\Q{21}{\dfrac{2x-1}{3}-\dfrac{x+1}{2}=1}\par\vspace{3mm}
\setlength{\ansh}{9.5mm}\setlength{\ansdrop}{6mm}%
\begin{center}\begin{tabular}{|p{\anscolw}|p{\anscolw}|p{\anscolw}|}\hline
\anscell{1} & \anscell{2} & \anscell{3} \\\hline
\anscell{4} & \anscell{5} & \anscell{6} \\\hline
\anscell{7} & \anscell{8} & \anscell{9} \\\hline
\anscell{10} & \anscell{11} & \anscell{12} \\\hline
\anscell{13} & \anscell{14} & \anscell{15} \\\hline
\anscell{16} & \anscell{17} & \anscell{18} \\\hline
\anscell{19} & \anscell{20} & \anscell{21} \\\hline
\end{tabular}\end{center}

\end{document}
